\uuid{1qhG}
\titre{Changements de variables}
\auteur{hanani ; mzari}
\organisation{exo7}
\licence{CC BY-NC-SA 4.0}
\datecreate{2020-01-20}
\niveau{L1}
\module{Analyse}
\chapitre{Calcul d'intégrales}
\sousChapitre{Autre}
\difficulte{2}
\theme{Calculs d'intégrales, changements de variables}
\exo7qcmid{lille-2.integrales.008}

\contenu{
\begin{question}
\qclassification{127.99}
\qtags{motcle=changements de variables}
Parmi les affirmations suivantes, cocher celles qui sont vraies :
\begin{answers}
\bad{Le changement de variable $t=\ln x$ donne $\displaystyle \int _1^{\mathrm{e}}\frac{\ln x}{x}\mathrm{d}x=\int _1^{\mathrm{e}}t\, \mathrm{d}t=\frac{\mathrm{e}^2-1}{2}$.}
\bad{Le changement de variable $t=1-x^2$ donne $\displaystyle \int _0^{1}2x\mathrm{e}^{1-x^2}\mathrm{d}x=-\int _0^{1}\mathrm{e}^{t}\mathrm{d}t$.}
\good{Le changement de variable $t=1+\mathrm{e}^x$ donne $\displaystyle \int _0^{\ln 3}\frac{\mathrm{e}^x}{1+\mathrm{e}^x}\mathrm{d}x=\ln 2$.}
\good{Pour tout réel $a>0$, on a : $\displaystyle \int _{-a}^{a}\frac{\sin x\, \mathrm{d}x}{1+\cos ^2x}=0$.}
\end{answers}
\vskip3mm
\begin{explanations}
Le changement de variable $t=\ln x$ donne : $\displaystyle \int _1^{\mathrm{e}}\frac{\ln x}{x}\mathrm{d}x=\int _0^1t\, \mathrm{d}t=\frac{1}{2}$. Avec $t=1-x^2$ dans la seconde intégrale, on obtient : $\displaystyle \int _0^{1}2x\mathrm{e}^{1-x^2}\mathrm{d}x=-\int _{1}^0\mathrm{e}^{t}\mathrm{d}t=\mathrm{e}-1$. Ensuite, avec $t=1+\mathrm{e}^x$ dans la troisième intégrale, on obtient : $\displaystyle \int _0^{\ln 3}\frac{\mathrm{e}^x}{1+\mathrm{e}^x}\mathrm{d}x=\int _2^4\frac{\mathrm{d}t}{t}=\ln 2$. Enfin, la fonction $\displaystyle x\mapsto \frac{\sin x}{1+\cos ^2x}$ est impaire. Donc, pour tout $a>0$ : $\displaystyle \int _{-a}^{a}\frac{\sin x\, \mathrm{d}x}{1+\cos ^2x}=0$.
\end{explanations}
\end{question}
}
